\documentclass[12pt,a4paper,oneside]{article}
\usepackage[brazil]{babel}
\usepackage[utf8]{inputenc}
\usepackage{olymp}

\setcounter{problem}{2}

\begin{document}

\begin{problem}{Coeficiente de Rendimento Acumulado}{cra}{2}
 
Na UFV, o Coeficiente de Rendimento Acumulado é um índice que mede o desempenho acadêmico do estudante e é obtido pela média ponderada das notas obtidas nas disciplinas cursadas, considerando como peso o número de créditos das respectivas disciplinas. É calculado pela fórmula:

\vspace{-6pt}
\begin{equation}
    CRA = \frac{\sum{(NF \times C})}{\sum{C}}\notag
\end{equation}

\vspace{-6pt}
em que $\sum{}$ é o somatório, NF é a nota final da disciplina e C é o número de créditos da disciplina.

Por exemplo, um estudante que tenha cursado 4 disciplinas, com os resultados da tabela abaixo,

\begin{center}
\begin{tabular}{|l|cccc|}
    \hline
    nota final  & 70 & 60 & 55 & 80 \\
    \hline
    créditos    & 4  & 6  &  4 &  2 \\
    \hline
\end{tabular}
\end{center}

terá seu CRA calculado da seguinte forma:

\begin{equation}
    CRA = \frac{70\times4 + 60\times6 + 55\times4 + 80\times 2}{4+6+4+2} = 63,75\notag
\end{equation}

Faça um programa para calcular o CRA, dadas as notas finais e créditos das disciplinas cursadas.

%\Notes
%
%Observações, se tiver.

\InputFile

A entrada contém três linhas. A primeira contém um único número natural $N$, indicando o número de disciplinas cursadas pelo estudante. A segunda linha contém $N$ valores inteiros $NF_i$, indicando a nota final obtida nas disciplinas. A terceira linha contém $N$ valores inteiros $C_i$ indicando o número de créditos das respectivas disciplinas.

Restrições: $1 \leq N \leq 50$, $0 \leq NF_i \leq 100$, $1 \leq C_i \leq 10$.


\OutputFile

Seu programa deve gerar apenas uma linha de saída, contendo o valor do CRA com 1 casa decimal.

% \Notes
% Preste atenção aos tipos das variáveis, pois $F(90) \approx 2^{61}$.

\Example

\begin{example}
\exmp{
4
70 60 55 80
4 6 4 2
}{
63.8
}%
\end{example}

\begin{example}
\exmp{
1
59
2
}{
59.0
}%
\end{example}

\begin{example}
\exmp{
6
100 100 60 60 60 39
6 5 4 3 2 1
}{
80.0
}%
\end{example}

%Comentários sobre o exemplo, se necessário.

\end{problem}

\end{document}
